On a Web increasingly read by generative agents, publishing a face and consenting
to model customization have become technically indistinguishable. Existing
anti-customization methods are evaluated on pristine files, although an agent
may acquire a photograph after a single ordinary Web operation such as CDN
transcoding, thumbnail rendering, screenshot cropping, restoration,
regeneration, or purification. We introduce \emph{ScrapeBench}, an acquisition
contract that evaluates each such operation independently from the same upload,
with a matched clean branch and held-out identity references. Our feasibility
pilot fixes four identities, with four gallery and two reference images per
identity, before inspecting defense outcomes. Completed comparisons cover
clean, GAP-Diff, MetaCloak, GoodAC, and ARCP-Web under InstantID and IP-Adapter-FaceID.
ARCP-Web allocates a bounded perturbation by texture and explicitly
optimizes the currently weakest of seven independent acquisition conditions.
We report conditional and end-to-end identity similarity, face-detection
failure, image quality, and paired identity-bootstrap intervals. Among 11
reader--layer conditions passing the clean validity gate, ARCP-Web obtains the
largest paired protection gap in nine, while respecting an 8/255 pixel budget
and attaining 37.80 dB PSNR. The small pilot is an effect-direction and systems
check, not a population estimate; a pre-audited 40-identity pool is used only
after its gates pass.
